% concept01.tex
\section{Ratios}

A ratio is about comparison. It tells us how much of one thing there is compared to
another. Ratios are everywhere: the number of players on a team, the ingredients in
a recipe, or the colors mixed for paint. A ratio doesn't just give numbers—it shows
balance.

If you have 2 apples and 3 oranges, the ratio of apples to oranges is 2 to 3. You
can write it \(2:3\), \(\frac{2}{3}\), or ``2 to 3.'' They all mean the same
thing.

Ratios also show part-to-whole relationships. If a class has 12 boys and 8 girls,
the ratio of boys to girls is \(12:8\), which simplifies to \(3:2\). But the ratio
of girls to total students (20) is \(8:20\), or \(2:5\). Ratios let us see
relationships clearly.

\blockquote{A ratio compares two numbers or quantities. It tells us how much of one
thing there is compared to another.}

\subsection{What is a Ratio?}

A \textbf{ratio} compares two numbers or quantities. It tells us how much of one
thing there is compared to another. Ratios help us describe relationships and
patterns, and they appear in recipes, sports, and even classrooms.

\textbf{Examples:}
\begin{itemize}
  \item \textbf{Part-to-Part:} apples to oranges \(\rightarrow\) \(2:3\).
  \item \textbf{Part-to-Whole:} apples to total fruit \(\rightarrow\) \(2:5\).
\end{itemize}

\textbf{Ways to write ratios:}
\begin{itemize}
  \item With a colon: \(\mathbf{2:3}\)
  \item As a fraction: \(\mathbf{\frac{2}{3}}\)
  \item In words: \(\mathbf{2\ \text{to}\ 3}\)
\end{itemize}
All three ways mean the same thing.

\textbf{Part-to-Part and Part-to-Whole:}
\begin{itemize}
  \item \textbf{Part-to-Part:} Compares two individual parts (e.g., apples to oranges \(\rightarrow\) \(2:3\)).
  \item \textbf{Part-to-Whole:} Compares one part to the total (e.g., apples to total fruits \(\rightarrow\) \(2:5\)).
\end{itemize}

\textbf{Example:} If a class has 12 boys and 8 girls:
\begin{itemize}
  \item The ratio of boys to girls is \(\mathbf{12:8}\) \(\rightarrow\) simplified to \(\mathbf{3:2}\).
  \item The ratio of girls to total students (20) is \(\mathbf{8:20}\) \(\rightarrow\) simplified to \(\mathbf{2:5}\).
\end{itemize}

Ratios simplify comparisons and help us see patterns.

\subsection{Did You Know?}

\begin{itemize}
  \item Ratios are used to mix paint colors, create scale models, adjust ingredients in recipes, and describe odds in games.
  \item Ratios are common in recipes, sports, and many real-world mixtures.
\end{itemize}

\subsection{Tips for Students}

\begin{itemize}
  \item Keep the order of comparison the same—read carefully.
  \item Simplify ratios when possible, just like you simplify fractions.
  \item Think visually—use pictures or charts if it helps you understand.
\end{itemize}


% concept02.tex
\section{Rates}

A rate is a special kind of ratio. It compares different units. We use rates to
measure speed, prices, and even how fast we type. A car's speed is miles per hour.
A store's price is dollars per item. A printer's speed is pages per minute.

If you drive 100 kilometers in 2 hours, the rate is
\[
\frac{100}{2}=50
\]
kilometers per hour. Rates make comparisons fair because they put everything on the
same scale.

\blockquote{A rate is a special type of ratio that compares two different kinds of
units.}

\subsection{What is a Rate?}

A \textbf{rate} is a special type of ratio that compares two different kinds of
units. For example, it can compare miles to hours, dollars to items, or words to
minutes.

\textbf{Examples:}
\begin{itemize}
  \item \(60\) pages in \(3\) minutes \(\rightarrow\) \(20\) pages per minute.
  \item \(\$200\) earned in \(8\) hours \(\rightarrow\) \(\$25\) per hour.
  \item If you drive \(100\) kilometers in \(2\) hours, the rate is \(\displaystyle \frac{100}{2} = \mathbf{50}\ \text{km/h}\).
\end{itemize}

\textbf{Common rates:}
\begin{itemize}
  \item \textbf{Speed:} miles per hour (mph), kilometers per hour (km/h).
  \item \textbf{Prices:} dollars per item (\$/item).
  \item \textbf{Typing:} words per minute (wpm).
\end{itemize}

\textbf{Everyday Examples:}
\begin{itemize}
  \item A printer prints \(60\) pages in \(3\) minutes \(\rightarrow\) Rate: \(\mathbf{20}\) pages per minute.
  \item A worker earns \(\$200\) for \(8\) hours \(\rightarrow\) Rate: \(\mathbf{\$25}\) per hour.
\end{itemize}

Rates help us measure and compare in real life.

\subsection{Did You Know?}

Rates are everywhere—in recipes, gas mileage (MPG), data usage on your phone, and
how fast athletes run.

\subsection{Tips for Students}

\begin{itemize}
  \item Always include the units (like miles, hours, or dollars).
  \item To find a rate, divide the first number by the second.
  \item Use rates to compare which deal is better when shopping.
\end{itemize}


% concept03.tex
\section{Unit Rate}

Sometimes we want to know the amount ``per one.'' That's called a unit rate. If 3
notebooks cost \$6, then each notebook costs \$2. Unit rates are powerful because
they help us compare fairly—per item, per mile, per hour.

When grocery shopping, you might see two boxes of cereal: one \$4 for 8 ounces,
another \$6 for 12 ounces. Unit rate shows which is the better deal.

\blockquote{A unit rate is a rate where the second number is \(1\). It tells you
how much of something there is per one unit of something else.}

\subsection{What is a Unit Rate?}

A \textbf{unit rate} is a rate where the second number is \(1\). It tells you how
much of something there is \emph{per one unit} of something else.

\textbf{Examples:}
\begin{itemize}
  \item \(120\) miles in \(4\) hours \(\rightarrow\) \(30\) miles per hour.
  \item \(\$4\) for \(8\) apples \(\rightarrow\) \(\$0.50\) per apple.
  \item If \(3\) notebooks cost \(\$6\), the unit rate is \(\displaystyle \frac{6}{3} = \mathbf{\$2}\) per notebook.
\end{itemize}

\textbf{How to find a unit rate:}
\begin{enumerate}
  \item Divide the first number by the second.
  \item Make sure the second number becomes \(1\).
\end{enumerate}

\textbf{Example:}
\[
120\ \text{miles in }4\ \text{hours} \rightarrow 120\div 4 = 30\ \text{miles per hour}.
\]
\[
8\ \text{apples for }\$4 \rightarrow \$4\div 8 = \$0.50\ \text{per apple}.
\]

\textbf{Real-life uses:}
\begin{itemize}
  \item Grocery stores list unit prices on shelf tags to help shoppers compare.
  \item Runners track their pace (minutes per mile) to measure speed.
  \item Restaurants calculate price per serving to set menu prices.
\end{itemize}

Unit rates appear in budgets, shopping, sports, and cooking.

\subsection{Did You Know?}

Unit prices in stores are unit rates—they show cost per ounce or per item.

\subsection{Tips for Students}

\begin{itemize}
  \item Make sure the second number becomes \(1\) (that's what ``unit'' means).
  \item Use real-world examples when practicing—like prices, distance, or time.
  \item Always include the unit in your answer (like ``per hour'' or ``per item'').
\end{itemize}


% concept04.tex
\section{Proportions}

Proportions show equality between two ratios. They tell us that two comparisons are
balanced. If a recipe for 2 people uses 1 cup of rice, then a recipe for 4 people
uses 2 cups. That's a proportion: the ratio stays the same.

We solve proportions by cross multiplying. If one ratio equals another, their cross
products match. This makes it easy to find missing values.

\blockquote{A proportion is a statement that shows two ratios or rates are equal.
It helps us solve problems where quantities grow or shrink at the same rate.}

\subsection{What is a Proportion?}

A \textbf{proportion} is a statement that shows two ratios or rates are \emph{equal}.
It helps us solve problems where quantities grow or shrink at the same rate.

\textbf{Example:}
\[
\frac{2}{3} = \frac{4}{6}
\]
If \(\frac{2}{3}=\frac{4}{6}\), this is a proportion because both sides are equal.

\textbf{How to solve proportions:} Use \emph{cross multiplication}:
\[
\frac{2}{3}=\frac{4}{6}
\quad\Longrightarrow\quad
2\times 6 = 3\times 4
\quad\Longrightarrow\quad
12 = 12.
\]
This means it is a true proportion.

You can also use proportions to find unknown numbers. For example,
\[
\frac{x}{5}=\frac{6}{10}
\quad\Longrightarrow\quad
10x = 30 \quad\Longrightarrow\quad x = 3.
\]

\textbf{Everyday examples:}
\begin{itemize}
  \item A map says \(1\) inch \(=\) \(50\) miles \(\rightarrow\) use this ratio to calculate real distances.
  \item A recipe that serves \(2\) people can be doubled to serve \(4\) using proportions.
\end{itemize}

Proportions keep comparisons fair and balanced.

\subsection{Did You Know?}

Proportions are used in architecture, photography, and fashion design to maintain
balance and scale.

\subsection{Tips for Students}

\begin{itemize}
  \item Always simplify ratios before comparing.
  \item Use cross multiplication to solve quickly.
  \item Write word problems as proportions to organize your thinking.
\end{itemize}


% concept05.tex
\section{Percents}

A percent means ``per hundred.'' It is another way to compare numbers. \(50\%\) is
half, \(25\%\) is a quarter, and \(100\%\) is the whole. Percents are part of
everyday life—on price tags, test scores, and bank interest.

Converting between fractions, decimals, and percents is simple once you know the
rules. Fractions turn into percents by multiplying by \(100\). Decimals turn into
percents by moving the decimal two places to the right. Percents turn back into
decimals by moving the decimal two places to the left.

\blockquote{A percent shows part of a whole, out of \(100\). It is a number that
means ``per 100.''}

\subsection{What is a Percent?}

A \textbf{percent} shows part of a whole, out of \(100\). It is a number that means
``\emph{per 100}.'' It shows a part out of a whole divided into \(100\) equal parts.

\textbf{Examples:}
\begin{itemize}
  \item \(50\%\) means \(50\) out of \(100\) \(\rightarrow\) \(\frac{1}{2}\).
  \item \(25\%\) means \(25\) out of \(100\) \(\rightarrow\) \(\frac{1}{4}\).
  \item \(100\%\) means all of something.
\end{itemize}

\subsection{Conversions}

\begin{itemize}
  \item Fraction \(\rightarrow\) Percent: \(\displaystyle \frac{1}{4}=0.25=25\%\).
  \item Decimal \(\rightarrow\) Percent: \(0.6 = 60\%\).
  \item Percent \(\rightarrow\) Decimal: \(45\% = 0.45\).
\end{itemize}

\textbf{How to convert:}
\begin{itemize}
  \item \textbf{Fraction to Percent:} multiply by \(100\).\\
    \(\displaystyle \frac{1}{4}=0.25 \Rightarrow 25\%\).
  \item \textbf{Decimal to Percent:} move the decimal point two places to the right.\\
    \(0.6 \Rightarrow 60\%\).
  \item \textbf{Percent to Decimal:} move the decimal point two places to the left.\\
    \(45\% \Rightarrow 0.45\).
\end{itemize}

\subsection{Real-Life Examples}

\begin{itemize}
  \item \textbf{Shopping:} \(30\%\) off a jacket means you save \(30\) out of every \(100\) dollars.
  \item \textbf{Grades:} A score of \(85\%\) means you got \(85\) out of \(100\) points.
  \item \textbf{Interest Rates:} Banks use percents to calculate how much money you earn or owe.
\end{itemize}

Percents help us see proportions quickly and clearly.

\subsection{Did You Know?}

The percent sign (\%) was first used in Europe hundreds of years ago. Today,
percents are used globally in business, science, and sports statistics.

\subsection{Tips for Students}

\begin{itemize}
  \item Think of \(\%\) as ``out of \(100\).''
  \item Use visual aids like \(100\)-grids or pie charts to understand percent.
  \item Practice converting between fractions, decimals, and percents until it feels natural.
\end{itemize}


% ending.tex
\section{Tip for Students}

Ratios, rates, proportions, and percents are not just classroom topics. They
describe how things relate and change. Learn to simplify ratios, find unit rates,
solve proportions, and convert percents. These tools will help you shop smarter,
plan better, and understand the world with clarity.
